\documentclass[12pt]{article}

% --------------------
% Packages
% --------------------
\usepackage[margin=1in]{geometry}     % margins
\usepackage{setspace}                 % line spacing
\usepackage{parskip}                  % no indent, space between paragraphs
\usepackage{graphicx}                 % images
\usepackage{amsmath, amssymb}         % math
\usepackage{booktabs}                 % tables
\usepackage{hyperref}                 % links
\usepackage{natbib}                   % citations
\usepackage{titlesec}
\usepackage{float}
\usepackage{placeins}
\usepackage{lineno}
\linenumbers
% section formatting

% --------------------
% Formatting
% --------------------
% \onehalfspacing                      % 1.5 spacing 
\doublespacing

\setlength{\parindent}{0pt}          % no paragraph indent
\setlength{\parskip}{0.5em}          % space between paragraphs


\titleformat{\section}{\large\bfseries}{\thesection}{1em}{}
\titleformat{\subsection}{\normalsize\bfseries}{\thesubsection}{1em}{}

% --------------------
%
% --------------------
\title{Climate-Driven Species Redistribution Increases the Vulnerability of TURFs}
%I wonder if i should mention "habitat suitability" over "spp redistribution" for accuracy, but im still unsure! I'll do my title last 
%Climate-driven species redistribution makes TURFs vulnerable to climate change
%Species redistribution models showcase TURFs are vulnerable to climate change
\author{Sabirah M. Carfagna}
\date{\today}

% --------------------
% Document
% --------------------
\begin{document}

\maketitle

\section{Introduction}

\subsection{Territorial Use Rights for Fisheries}

Territorial Use Rights for Fisheries (TURFs) are spatially defined management systems that allocate exclusive access rights over marine resources within a defined geographic area to a well defined group of fishers \citep{McCayEtAl2014, McCay2017}. In doing so, TURFs reduce overexploitation associated with open-access dynamics \citep{McCayEtAl2014, McCay2017}. Because the benefits of conservation, harvest restraint, and local monitoring accrue to the same users who depend on the resource over time, TURFs encourage long-term stewardship and sustainable resource use \citep{Gelcich2015, Viana2019}. Within these systems, resource users are often responsible for local monitoring, surveillance, enforcement, and day-to-day management of fishing activities within their territories, while broader regulatory frameworks, including TURF allocation, harvest regulations, quotas, and other fisheries management measures, are established by government agencies \citep{McCay2017, McCayEtAl2014, Ojea2017, Gelcich2015}.

TURF-like fisheries management systems have been implemented in many places around the world, including Mexico, Chile, Belize, Japan, Brazil, and the Philippines, among many others \citep{Aceves-Bueno2017-qb, Wade2019, Matsuda2010, Siar1992}. Despite differences in their legal and institutional structures, these systems share the common feature of granting exclusive access rights to defined fishing areas and resources. Evidence from Chile, where TURFs take the form of Management and Exploitation Areas for Benthic Resources (MEABRs), suggests that well-managed TURFs can produce substantial ecological and social benefits \citep{Gelcich2008, Gelcich2010, Gelcich2015}. These benefits include higher species richness, biomass, and density relative to adjacent open-access areas, recovery of benthic resources such as loco (\textit{Concholepas concholepas}), and strengthened stewardship and collective action among fishing organizations \citep{Gelcich2008, Gelcich2010, Gelcich2015}.

A key condition underlying the effectiveness of TURFs is the alignment between management systems and the ecological characteristics of the resources they are intended to govern, an idea commonly discussed within the broader framework of social-ecological fit \citep{PoonBonzon2013, AcevesBueno2018, Nenadovic2026}. One important aspect of this alignment is the commensurability between the spatial scales of management institutions and the ecological processes operating within the resource system, enabling resource users and management institutions to respond effectively to ecological change \citep{McCay2017, Nenadovic2026}. In practice, this means that the ecological characteristics of the resource, including species movement, habitat distribution, connectivity, and recruitment processes, must remain aligned with TURF boundaries \citep{McCay2017, PoonBonzon2013, AcevesBueno2018}.

This relationship can be understood through the interaction of biological and social functional units. Biological functional units correspond to the geographic range of self-sustaining stocks or sub-stocks and may encompass spawning grounds, nursery habitats, and migratory routes, while social functional units consist of the groups of resource users capable of organizing and participating in fishery management \citep{PoonBonzon2013}. Species with localized biological functional units, such as those with limited movement, are generally more compatible with TURF systems because the ecological processes supporting those species occur at spatial scales that can be more readily influenced through local management actions \citep{PoonBonzon2013, Nenadovic2026}. This fit, however, may be threatened by climate change. Climate change may drive changes in species distributions that modify the spatial correspondence between biological and social functional units within existing TURF boundaries.

\subsection{Species Distributions and Climate Change}

Species distributions are determined by a combination of abiotic and biotic conditions that influence where species can survive, reproduce, and maintain self-sustaining populations. Abiotic conditions include environmental characteristics such as temperature, oxygen, salinity, and depth, whereas biotic factors include ecological relationships such as competition, predation, and food availability \citep{MeloMerino2020}. In marine systems, abiotic conditions are particularly important because species can only persist within a finite range of environmental conditions that support essential physiological processes \citep{Portner2007}. As a result, variation in factors such as temperature, oxygen availability, and salinity can directly influence species performance and ultimately determine where populations are able to occur. This is particularly true for marine ectotherms whose physiology is strongly influenced by surrounding environmental conditions \citep{Nguyen2011, Portner2007}.

Climate change is modifying marine environments through a range of physical and chemical changes, including ocean warming, declining dissolved oxygen concentrations, changes in salinity patterns, ocean acidification, and altered ocean productivity \citep{Portner2007, Pecl2017}. Many of these changes are linked to the accumulation of anthropogenic greenhouse gases, which alter ocean circulation, stratification, and other large-scale oceanographic processes \citep{Portner2007}. Climate change is also altering salinity, nutrient availability, and seawater chemistry through changes in precipitation, runoff, circulation, and atmospheric carbon dioxide concentrations \citep{Portner2007}. Because these environmental variables influence the conditions under which species can survive, grow, reproduce, and maintain viable populations, climate-driven changes in ocean conditions can alter the location, extent, and quality of suitable habitat through time \citep{Portner2007, MeloMerino2020}. Species are therefore expected to respond to climate change by tracking shifts in suitable habitat, resulting in changes in geographic distributions and species ranges \citep{Pecl2017}.

A growing body of evidence indicates that marine species are already responding to climate-driven environmental change through shifts in their geographic distributions. Marine organisms have expanded, contracted, or relocated their ranges at average rates approaching 70 km per decade, although responses vary substantially among species and regions \citep{Pecl2017, Pinsky2018}. Although poleward and deeper range shifts are among the most commonly observed responses to ocean warming, species have exhibited considerable variation in the direction and magnitude of distributional change \citep{Pecl2017, Pinsky2013}. This variation in species redistributions has been linked to differences in climate velocity, which describes the rate and direction at which climatic conditions shift across space \citep{Pinsky2013}. Because species redistribute by tracking shifts in suitable habitat, projected changes in habitat suitability can provide a useful basis for anticipating potential changes in species distributions under future climate scenarios. As climate-driven redistributions continue, species are increasingly crossing ecological and political boundaries, with projections suggesting that 23-35\% of Exclusive Economic Zones (EEZs) may contain new transboundary stocks by the end of the century \citep{Pinsky2018}.

These climate-driven redistributions present a challenge for area-based management systems because management boundaries remain fixed while the distribution of the resources they are intended to govern may shift through time. This challenge has received increasing attention at larger governance scales, particularly in the context of transboundary stocks that cross Exclusive Economic Zone (EEZ) boundaries and require coordination among jurisdictions \citep{Pinsky2018, Pecl2017}. Because TURFs are geographically fixed management areas nested within EEZs, it is reasonable to expect that similar redistributions may also affect the spatial relationship between target species and TURF boundaries. As species and suitable habitat shift relative to fixed management areas, the ecological processes occurring within those areas may no longer correspond to those present when management systems were originally designed, potentially weakening the social-ecological fit and stewardship incentives upon which TURF effectiveness depends \citep{McCay2017, PoonBonzon2013, AcevesBueno2018, Nenadovic2026}. Understanding climate-driven changes in suitable habitat within TURFs is therefore essential for evaluating the vulnerability of these spatially fixed management systems.

\subsection{Objectives and Hypotheses}

Despite extensive work examining climate-driven species redistributions under future climate scenarios, most studies have focused on broad spatial scales, such as Exclusive Economic Zones (EEZs) and regional fisheries systems \citep{Pinsky2018, Pecl2017, Cheung2010}. These systems are mostly relevant to industrial and large-scale fisheries. In contrast, comparatively little attention has been given to how climate-driven changes in habitat suitability may affect smaller-scale territorial management systems such as TURFs, which are important for sustaining small-scale fisheries in tropical and developing countries. To our knowledge, this is the first study to quantify projected changes in species-specific habitat suitability within TURFs under future climate scenarios. This knowledge gap is particularly important because many TURF systems occur in tropical and subtropical regions, where species may be especially sensitive to climate change because they already occupy environmental conditions close to their physiological tolerances \citep{Portner2007, Nguyen2011, Pecl2017}. Consistent with this pattern, projections suggest that many tropical fisheries may experience declines in habitat suitability and fisheries productivity under climate change, while higher-latitude regions may experience gains as suitable environmental conditions shift poleward \citep{Cheung2010}. Similarly, climate-driven species redistributions are expected to occur through contractions at the warmer, equatorward portions of species distributions and expansions toward higher latitudes as suitable environmental conditions shift poleward \citep{FredstonHermann2020}.

Mexico provides a useful empirical system for evaluating how climate-driven changes in habitat suitability may affect spatially fixed fisheries management systems. The Mexican TURF system comprises more than 200 TURFs distributed across nine coastal states bordering the Pacific Ocean, Gulf of California, Gulf of Mexico, and Mexican Caribbean (\citet{AcevesBueno2023}; \autoref{fig:mexturfs}). Most Mexican TURFs are granted to fishing cooperatives and are issued for renewable 20-year periods, contingent on continued productivity and responsible resource management \citep{McCay2017, McCayEtAl2014, AcevesBueno2023}. Across the system, TURFs have been issued for 26 taxa. With the exception of finfish, all taxa are benthic, including shrimp (\textit{Penaeidae}), lobster (\textit{Palinuridae}), abalone (\textit{Haliotidae}), algae, and other marine resources \citep{AcevesBueno2023}. TURF size varies substantially, ranging from less than 1 km$^2$ to more than 9,000 km$^2$ \citep{AcevesBueno2023}.

\begin{figure}[tbh!]
\centering
\includegraphics[width=0.9\linewidth]{sab_mexplot.png}
\vspace{-8pt}
\caption{Spatial distribution of Mexican Territorial Use Rights for Fisheries (TURFs) included in this study. Grey polygons represent TURF boundaries compiled from \citet{AcevesBueno2023}.}
\label{fig:mexturfs}
\end{figure}


\FloatBarrier
Using Mexico as an empirical study system, I propose to examine how climate change may affect the persistence of target species within TURFs. To do so, I use current and future habitat suitability projections from AquaX, a marine species distribution modeling framework that estimates the relative suitability of environmental conditions for individual species using a Habitat Suitability Index (HSI). Because climate-driven species redistribution is expected to occur through shifts in suitable habitat, projected habitat suitability is used as a proxy for current and future species distributions. These data, however, do not allow us to track abundance or biomass or the respective populations. Accordingly, I will evaluate changes in mean habitat suitability, percent suitable habitat, and target species richness within TURFs to assess how climate change may alter the quantity and quality of suitable habitat and number of target species contained within existing TURF boundaries across three shared socioeconomic pathways (SSPs). These metrics, defined in detail in \autoref{sec:HS}, are used to test the following hypotheses:

\paragraph{$H_1$ (change through time):} Mean habitat suitability, percent suitable habitat, and target species richness within Mexican TURFs will \textit{decline} from the current conditions to mid- and end-century projections.

\paragraph{$H_2$ (change across climate scenarios):} Declines in mean habitat suitability, percent suitable habitat, and target species richness within Mexican TURFs will be \textit{larger} under SSP5–8.5 (most severe) than under SSP1–2.6 (significant mitigation).

\paragraph{$H_3$ (spatial vulnerability):}
Losses in mean habitat suitability, percent suitable habitat, and number of target species will be \textit{larger} in lower-latitude TURFs than in higher-latitude TURFs.

\section{Methods}

\subsection{Data Sources}

\subsubsection{Geospatial TURF Dataset}

Spatial data on Mexican TURFs are available from the nationwide concession database compiled by \citet{AcevesBueno2023}, which maps the geographic boundaries and institutional characteristics of each TURF across Mexico. The database was developed using official records with spatial coordinates obtained through the Mexican transparency system (INFOMEX) and CONAPESCA, which were subsequently digitized and mapped in GIS software to create polygon shapefiles for each TURF \citep{AcevesBueno2023}. The final dataset contains 214 polygons distributed across nine Mexican coastal states, with approximately 87\% located along the Pacific coast \citep{AcevesBueno2023}. Most TURF holders (97.1\%) possess a single TURF title, and the majority of TURFs (89.8\%) consist of a single polygon. Nearly all TURFs are held by cooperatives, with only three private entities identified in the database. The dataset additionally contains information on TURF holder, targeted taxa, TURF geometry, and TURF location \citep{AcevesBueno2023}. These data were complemented with additional TURF boundaries for TURFS along the Yucatan Peninsula, between the Gulf of Mexico and Caribbean Sea.

Approximately 66\% of TURFs in the Mexican TURF system are located in estuarine environments \citep{AcevesBueno2023}. These TURFs are concentrated in lagoon and estuarine systems along the Pacific and Gulf coasts, particularly in the states of Sinaloa, Sonora, Chiapas, and Veracruz. Because large-scale marine habitat suitability projections may not adequately resolve environmental conditions within estuarine, lagoon, and other enclosed coastal systems, this study excludes estuarine TURFs and focuses exclusively on marine TURFs.

\subsubsection{Habitat Suitability Projections}

Habitat suitability projections for target species will be obtained from AquaX, an updated marine species distribution modeling framework developed by \citet{Reygondeau2026}. AquaX begins with species occurrence records compiled from AquaMaps, the Ocean Biodiversity Information System (OBIS), and the Global Biodiversity Information Facility (GBIF), with taxonomy standardized using the World Register of Marine Species (WoRMS). Biological and biogeographical information from FishBase and SeaLifeBase, including habitat preferences and depth ranges, is then used to refine these records and constrain predictions to biologically plausible habitats. The resulting verified species occurrences are related to environmental conditions, including temperature, dissolved oxygen, salinity, productivity, pH, seawater velocity, and depth, depending on species habitat type. These relationships are modeled using ten species distribution modeling algorithms. Individual models are evaluated using the Area Under the Curve (AUC), True Skill Statistic (TSS), and Critical Success Index (CSI), and models meeting the performance criterion are retained and combined into an ensemble prediction. The resulting ensemble estimates habitat suitability for each species as a Habitat Suitability Index (HSI), which represents the relative suitability of environmental conditions for that species. AquaX then applies these modeled species–environment relationships to contemporary and projected future environmental conditions to generate spatial habitat suitability projections under different climate scenarios \citep{Reygondeau2026}.

AquaX produces habitat suitability projections at 0.05° resolution, approximately 5 km at the equator, for a contemporary baseline period (2000–2010) and two future periods, mid-century (2050–2060) and end-century (2090–2100), under three Shared Socioeconomic Pathways: SSP1–2.6, SSP2–4.5, and SSP5–8.5 \citep{Reygondeau2026}.

AquaX expresses habitat suitability as a Habitat Suitability Index (HSI) ranging from 0 to 1. Suitability is highest (1.0) within a species' preferred environmental range and declines toward zero as environmental conditions approach thresholds considered unsuitable for survival \citep{Reygondeau2026}. Comparing habitat suitability at a given location across time periods can indicate whether environmental conditions for a species are projected to decline, persist, or improve. Because AquaX provides both contemporary and future habitat suitability projections, these data can be overlaid with TURF polygons to quantify changes in habitat suitability within TURFs over time.

\subsection{Data Processing}\label{sec:HS}

Species information associated with each TURF will be standardized prior to analysis to link TURF target species with AquaX habitat suitability projections. Target species in the Aceves-Bueno et al. (2023) database will be reviewed across TURFs. Species are recorded using a combination of Spanish common names, English common names, genus-level classifications, and scientific names, requiring species names to be standardized across the dataset. 
To ensure compatibility with AquaX projections, standardized scientific names and corresponding AphiaIDs from WoRMS will be assigned to each target species. In cases where species identifications are incomplete or ambiguous, additional information from TURF documentation and the Carta Nacional Pesquera will be used to verify the most likely species identity or taxonomic group represented by the TURF. The complete list of target taxa and corresponding AphiaIDs is provided in Supplementary \autoref{tab:target_taxa}.

\subsection{Spatial Analysis and Habitat Suitability Metrics}

Species-specific AquaX habitat suitability raster layers will be overlaid with TURF polygons to quantify projected changes in habitat suitability within management boundaries across climate scenarios and time periods. Three complementary metrics will be calculated: mean habitat suitability, representing habitat quality (\autoref{eq:mean_hsi}); percent suitable habitat, representing the extent of suitable habitat relative to TURF size (\autoref{eq:percent_suitable}); and target species richness, representing the number of target species that retain at least 10\% suitable habitat within a TURF (\autoref{eq:species_richness}). Mean habitat suitability and percent suitable habitat provide complementary information because the extent of suitable habitat within a TURF may change independently of the average quality of the habitat that remains suitable. For each species, suitable habitat will be identified using the species-specific HSI cutoff provided by AquaX. These cutoffs are determined during model calibration using the True Skill Statistic (TSS), which evaluates the model’s ability to correctly distinguish known species occurrences from absences. The cutoff identifies the HSI value that best separates suitable from unsuitable conditions for each species \citep{Reygondeau2026}. Because habitat suitability relationships differ among species, each species will be evaluated using its corresponding AquaX cutoff.


\subsubsection{Mean Habitat Suitability}

I define mean habitat suitability $\overline{HSI}_{ij}$ for species $j$ in TURF $i$ as:

\begin{equation}
    \overline{HSI}_{ij} = \frac{1}{n_{ij}} \sum_{k=1}^{n_{ij}} HSI_{ijk}
    \label{eq:mean_hsi}
\end{equation}

where $HSI_{ijk}$ represents the habitat suitability value of raster cell $k$, and $n_{ij}$ represents the number of raster cells within TURF $i$ with HSI values exceeding the species-specific AquaX cutoff for species $j$. By restricting calculations to suitable cells, this metric represents the average quality of suitable habitat for each target species within a TURF. Declines in mean HSI through time may indicate decreasing environmental suitability for target species within TURFs. Mean habitat suitability will be calculated for each TURF-species combination across the three climate scenarios and two time periods to quantify changes in the quality of suitable habitat within TURFs (See \autoref{sec:data_analysis}). 


\subsubsection{Percent Suitable Habitat}

I define percent suitable habitat (\%$S_{ij}$) for species $j$ within TURF $i$ as:

\begin{equation}
\%S_{ij} =
\frac{\sum_{k=1}^{N_i} I(HSI_{ijk}>c_j)}
{N_i}\times100
\label{eq:percent_suitable}
\end{equation}

where $N_i$ represents the total number of raster cells within TURF $i$, $c_j$ represents the species-specific AquaX HSI cutoff for species $j$, and $I(HSI_{ijk}>c_j)$ is an indicator function equal to 1 when habitat suitability exceeds the species-specific cutoff and 0 otherwise. This metric represents the proportion of TURF area containing environmentally suitable habitat for a target species. Declines in percent suitable habitat indicate loss of suitable habitat within TURF boundaries. Percent suitable habitat will be calculated for each TURF-species combination across the three climate scenarios and two time periods to quantify changes in the extent of suitable habitat within TURFs (See \autoref{sec:data_analysis}). 

\subsubsection{Species Richness}

I define species richness ($SR_i$) within TURF $i$ as:

\begin{equation}
    SR_i = \sum_j I(\%S_{ij}>\tau)
    \label{eq:species_richness}
\end{equation}

where $\tau = 10\%$ represents the minimum proportion of suitable habitat used to classify a target species as retained within a TURF, and $I(\%S_{ij}>\tau)$is an indicator function equal to 1 when the proportion of suitable habitat exceeds this threshold and 0 otherwise. This metric represents the number of target species retaining at least $10\%$ suitable habitat within a TURF. Because the $10\%$ threshold is an analytical choice used to translate continuous changes in suitable habitat into species richness, sensitivity analyses using alternative values of $\tau$ will be conducted to evaluate whether patterns in species richness are robust to this choice. Changes in species richness through time may indicate changes in the retention or loss of suitable habitat for target species within TURFs under future climate scenarios. Species richness will be calculated for each TURF across the three climate scenarios and two time periods (see \autoref{sec:data_analysis}).

\subsection{Data Analysis}\label{sec:data_analysis}

Temporal changes in habitat suitability metrics across climate scenarios and time periods will be visualized using comparative plots of mean habitat suitability, percent suitable habitat, and species richness across the contemporary baseline (2000-2010), mid-century (2050-2060), and end-century (2090-2100) projections under the Shared Socioeconomic Pathway scenarios SSP1-2.6, SSP2-4.5, and SSP5-8.5. Relationships between habitat quality and habitat extent will additionally be examined by plotting mean HSI against percent suitable habitat to evaluate different trajectories of habitat change within TURFs, including simultaneous declines in habitat quality and extent or shifts where habitat extent changes while habitat quality remains relatively stable.

The stated hypotheses will be formally evaluated using multiple linear regression models. Separate models will be fit for each of the three outcomes of interest outlined above.

To test $H_1$ (i.e. that mean habitat suitability, percent suitable habitat, and species richness will decline from present-day to future climate projections) changes in each outcome of interest relative to current conditions will be modeled as a function of projection period:

\begin{equation}
\Delta y_{ijts} = \beta_1 M_{ts} + \beta_2 E_{ts} + \varepsilon_{ijts}
\label{eq:h1_regression}
\end{equation}

where $\Delta y_{ijts} = y_{ijts} - y_{ij,\mathrm{current}}$ represents the change in outcome of interest $y$ for species $j$ in TURF $i$ under period $t$ and scenario $s$, relative to current conditions. The indicator variables $M_{ts}$ and $E_{ts}$ represent mid-century and end-of-century projections, respectively. The coefficients $\beta_1$ and $\beta_2$ represent the average change in the outcome of interest from present-day conditions to mid-century and end-of-century projections across all Shared Socioeconomic Pathway scenarios. The error term $\varepsilon_{ijts}$ represents variation in the outcome of interest not explained by projection period.

Under the null hypothesis, $\beta_1 = \beta_2 = 0$, indicating no change through time. Support for H1 would be indicated by negative coefficient estimates for both coefficients, with the magnitude of decline being greater by the end of the century than by mid-century, so $\beta_2 < \beta_1 < 0$.

To test H2, (i.e. larger declines under higher-emissions scenarios) changes in each outcome of interest will be modeled as a function of projection period, SSP scenario, and their interaction according to:


\begin{equation}
\begin{split}
\Delta y_{ijts} =
&\beta_1 M_{ts} +
\beta_2 E_{ts} +
\delta_1 \left(M_{ts} \times \text{SSP2-4.5}_{s}\right) +
\delta_2 \left(M_{ts} \times \text{SSP5-8.5}_{s}\right) + \\
&\delta_3 \left(E_{ts} \times \text{SSP2-4.5}_{s}\right) +
\delta_4 \left(E_{ts} \times \text{SSP5-8.5}_{s}\right) +
\varepsilon_{ijts}
\end{split}
\label{eq:h2_regression}
\end{equation}

where $M_{ts}$ and $E_{ts}$ represent mid-century and end-of-century projections, respectively, and SSP1-2.6 is treated as the reference scenario. The coefficients $\beta_1$ and $\beta_2$ therefore represent the average change in the outcome of interest under SSP1-2.6 at mid-century and end-of-century, respectively. The coefficients $\delta_1$ and $\delta_2$ represent the additional change associated with SSP2-4.5 and SSP5-8.5, respectively, relative to SSP1-2.6 at mid-century. Similarly, $\delta_3$ and $\delta_4$ represent the additional change associated with SSP2-4.5 and SSP5-8.5, respectively, relative to SSP1-2.6 at end-of-century.

Under the null hypothesis, $\delta_1 = \delta_2 = \delta_3 = \delta_4 = 0$, indicating that projected changes do not differ among SSP scenarios. Support for H2 would be indicated by increasingly negative SSP coefficients under higher-emissions scenarios, such that declines under SSP5-8.5 exceed those under SSP2-4.5, and declines under SSP2-4.5 exceed those under SSP1-2.6. Thus, support for H2 would be indicated when $\delta_2 < \delta_1 < 0$ at mid-century and $\delta_4 < \delta_3 < 0$ at end-of-century.

To test H3, which predicts that climate impacts will vary spatially across Mexican TURFs, changes in each outcome of interest will be modeled as a function of TURF latitude:

\begin{equation}
\Delta y_{ijts} = \alpha + \theta X_i + \varepsilon_{ijts}
\label{eq:h3_regression}
\end{equation}

where $\Delta y_{ijts}$ represents the change in mean habitat suitability, percent suitable habitat, or species richness for species $j$ in TURF $i$ under period $t$ and scenario $s$, relative to current conditions. The variable $X_i$ represents the latitude of the centroid of TURF $i$, and $\theta$ represents the effect of latitude on projected changes in the outcome of interest. The error term $\varepsilon_{ijts}$ represents variation in the outcome of interest not explained by latitude.

Under the null hypothesis, $\theta = 0$, indicating that projected changes do not vary with TURF latitude. Support for H3 would be indicated by a positive coefficient estimate ($\theta > 0$), suggesting that higher-latitude TURFs experience smaller declines in habitat suitability, suitable habitat extent, or species richness than lower-latitude TURFs.


 \section{Limitations and Delimitations}

This study is intentionally delimited to coastal marine TURFs within Mexico and excludes estuarine TURFs due to limitations associated with large-scale marine habitat suitability projections in enclosed coastal systems. Analyses are additionally restricted to species for which compatible AquaX habitat suitability projections are available. The study further focuses on projected changes in habitat suitability as a proxy for climate-driven species redistribution and does not directly evaluate observed abundance, biomass, fisheries yield, or socioeconomic adaptation within fishing communities.

Several methodological limitations should also be considered. AquaX habitat suitability projections represent modeled environmental suitability rather than confirmed species presence or abundance, and therefore changes in HSI should not be interpreted as direct measures of biomass or fisheries productivity \citep{Reygondeau2026}. Habitat suitability projections additionally inherit uncertainty associated with species occurrence records, environmental data resolution, future climate scenarios, and assumptions embedded within species distribution models \citep{Reygondeau2026}. Because AquaX projections are produced at a spatial resolution of approximately 0.05$^\circ$, fine-scale habitat heterogeneity and localized environmental processes may not be fully resolved, particularly within smaller TURFs. Finally, the 10\% threshold used to classify target species as retained within a TURF is an analytical choice that may influence estimates of species richness, and sensitivity analyses using alternative threshold values will therefore be used to evaluate the robustness of these results. 


\bibliographystyle{apalike}
\bibliography{references}

\clearpage
\appendix
\counterwithin{figure}{section}
\counterwithin{table}{section}
\renewcommand\thefigure{\thesection\arabic{figure}}
\renewcommand\thetable{\thesection\arabic{table}}

\section{Supplementary Figures and Tables}
\label{app:species}

%\input{species_table.tex}


\end{document}